\documentclass[12pt]{article}
\usepackage[utf8]{inputenc}
\usepackage[T1]{fontenc}
\usepackage{hyperref}
\title{Tunia Language Resources for K-12 Curriculum Development}
\author{Compiled by Claude Sonnet 5 (claude-sonnet-5) for Andr\'as Kornai}
\date{}
\begin{document}
\maketitle

\section{Introduction}
Tunia (ISO 639-3 tug; Glottocode tuni1251), also written Tounia or Tunya, is a Bua (Adamawa) language of Chad in the Atlantic-Congo family. It is spoken around Sarh in Moyen-Chari, Barh Kôh Department. Speaker figures are about 2,300 (1993 census) or 2,260. The Endangered Languages Project rates it as threatened, and its Perim dialect is reported extinct. The searches found only reference and demographic pages. No educational materials exist in the results.

\section{Primary Educational Resources}
None found.

\section{Linguistic and Reference Materials}
\begin{itemize}
\item Wikipedia overview of Tunia: \url{https://en.wikipedia.org/wiki/Tunia_language}
\item Wikipedia overview of the Bua languages: \url{https://en.wikipedia.org/wiki/Bua_languages}
\item Glottolog entry for Tunia: \url{https://glottolog.org/resource/languoid/id/tuni1251.bigmap.html}
\item Endangered Languages Project profile: \url{https://endangeredlanguages.com/node/90410}
\item Endangered Languages Project bibliography: \url{https://endangeredlanguages.com/lang/4586/bibliography}
\end{itemize}
The content of the bibliography was not opened, so whether it lists any usable description is unknown.

\section{Religious and Community Resources}
\begin{itemize}
\item Joshua Project profile of the Tunya people in Chad (9,100 people; this counts people, not speakers): \url{https://joshuaproject.net/people_groups/15645}
\end{itemize}

\section{Digital and Audio Resources}
None found.

\section{Curriculum Development Context}
There is no evidence of written materials, orthography or schooling in Tunia, and the language is likely taught only at home, if at all. Any curriculum would start from documentation work with speakers around Sarh, with French and Chadian Arabic as the likely languages of schooling.
\end{document}